\par\medskip
\hypertarget{native-086-article067}{}\pdfbookmark[2]{67.}{article-067}
\begin{center}67.\end{center}
Ein anderes Verfahren zur Lösung der Aufgabe des vorigen Artikels ergibt sich aus den Gleichungen\HistoricalNoteEight{Handschriftliche Bemerkung:]
\[
\begin{aligned}
\sin\alpha\cos\delta&=-\sin\epsilon\sin b+\cos\epsilon\cos b\sin l\\[.35em]
\cos\alpha\cos\delta&=\cos b\cos l\\[.35em]
\sin\delta&=\cos\epsilon\sin b+\sin\epsilon\cos b\sin l.
\end{aligned}
\]
}
\[
\begin{aligned}
\cos\epsilon\sin l&=\sin\epsilon\operatorname{tang}b+\cos l\operatorname{tang}\alpha\\[.4em]
\sin\delta&=\cos\epsilon\sin b+\sin\epsilon\cos b\sin l\\[.4em]
\cos b\cos l&=\cos\alpha\cos\delta
\end{aligned}
\]
Man bestimme den Hilfswinkel $\theta$ aus der Gleichung
\[
\operatorname{tang}\theta=\frac{\operatorname{tang}b}{\sin l},
\]
dann gilt
\[
\begin{aligned}
\operatorname{tang}\alpha&=\frac{\cos(\epsilon+\theta)\operatorname{tang}l}{\cos\theta}\\[.6em]
\operatorname{tang}\delta&=\sin\alpha\operatorname{tang}(\epsilon+\theta),
\end{aligned}
\]
\SourceLeaf{087}Zur Kontrolle der Rechnung kann man diesen Gleichungen hinzufügen
\[
\cos\delta=\frac{\cos b\cos l}{\cos\alpha}\quad\text{oder}\quad
\cos\delta=\frac{\cos(\epsilon+\theta)\cos b\sin l}{\cos\theta\sin\alpha}\cdot
\]
Die Mehrdeutigkeit bei der Bestimmung von $\alpha$ aus der zweiten Gleichung wird dadurch aufgehoben, dass $\cos\alpha$ und $\cos l$ dasselbe Vorzeichen haben müssen.

Dieses Verfahren ist weniger bequem, wenn außer $\alpha$ und $\delta$ auch $E$ gesucht ist. Die bequemste Formel zur Bestimmung dieses Winkels lautet dann $\cos E=\frac{\sin\epsilon\cos\alpha}{\cos b}=\frac{\sin\epsilon\cos l}{\cos\delta}\cdot$ Mit dieser Formel lässt sich $E$ aber nicht genau berechnen, wenn $\pm\cos E$ nur wenig von der Einheit abweicht. Außerdem bleibt offen, ob $E$ zwischen $0$ und $180^0$ oder zwischen $180^0$ und $360^0$ anzusetzen ist. Der erste Nachteil fällt selten ins Gewicht, zumal zur Berechnung der Differenzialquotienten keine höchste Genauigkeit im Wert von $E$ erforderlich ist. Die Mehrdeutigkeit lässt sich jedoch mithilfe der Gleichung $\cos b\cos\delta\sin E=\cos\epsilon-\sin b\sin\delta$ leicht beheben. Sie zeigt, dass $E$ zwischen $0$ und $180^0$ oder zwischen $180^0$ und $360^0$ anzusetzen ist, je nachdem, ob $\cos\epsilon$ größer oder kleiner als $\sin b\sin\delta$ ist. Offenbar ist diese Prüfung nicht einmal nötig, wenn einer der Winkel $b$, $\delta$ die Grenze $66^0\,32'$ nicht überschreitet. Dann ist nämlich $\sin E$ stets positiv. Übrigens kann dieselbe Gleichung im oben erwähnten Fall auch zur genaueren Bestimmung von $E$ verwendet werden, wenn sich die Mühe zu lohnen scheint.
